\documentclass[10pt]{article}
\usepackage{graphicx}
\setlength{\textwidth}{6.85in}
\setlength{\oddsidemargin}{-0.18in}
\setlength{\evensidemargin}{-0.18in}
\setlength{\textheight}{9.45in}
\setlength{\topmargin}{-0.65in}
\setlength{\parindent}{0pt}
\setlength{\parskip}{0.45em}
\setlength{\emergencystretch}{3em}
\title{\textbf{How We Classify Downwind and Upwind AC Area and Population}}
\author{Corrected East-referenced construction}
\date{August 2026}

\begin{document}
\maketitle

\section*{Executive summary}
For each focal grid and month, one rolling wind direction divides its assigned
assembly constituency (AC) into a downwind half and an upwind half. We then
construct two related but distinct indicators:
\[
\texttt{downup\_ac\_area}
=\mathbf{1}[A_{down}>A_{up}],
\qquad
\texttt{downup\_ac\_pop}
=\mathbf{1}[P_{down}>P_{up}].
\]
The area indicator compares exact polygon areas. The population indicator
compares the sum of 2010 population assigned to other small-grid centroids in
the same AC. Therefore, the indicators can differ when population is unevenly
distributed across the AC.

\section*{One angular convention throughout}
Every directional calculation is a Cartesian angle measured counterclockwise
from East:
\[
0^\circ=\mathrm{East},\quad 90^\circ=\mathrm{North},\quad
180^\circ=\mathrm{West},\quad 270^\circ=\mathrm{South}.
\]
The wind is $w=\mathrm{atan2}(v10,u10)$. For a comparison-grid centroid $j$
relative to focal grid $i$, the bearing is
\[
\beta_{ij}=\mathrm{atan2}(y_j-y_i,x_j-x_i).
\]
No step interprets either angle as a clockwise-from-North compass bearing.
The normalized calculation angle remains East-referenced; normalization only
changes its numeric range to $[0,360)$.

\section*{Real example used in the figure}
The reproducible script selected grid \textbf{77902} in
\textbf{Prithla} AC
(AC ID \textbf{262}),
Faridabad, Haryana, for
\textbf{2019-03}. The rolling wind direction is
\textbf{-57.820$^\circ$ from East}.

\begin{center}
\begin{tabular}{lrrr}
\hline
Measure & Downwind & Upwind & Indicator \\
\hline
Area (km$^2$) & 186.74 & 99.54 &
1 \\
Area share (\%) & 65.2 & 34.8 & \\
Population & 115438 & 152215 &
0 \\
Population share (\%) & 43.1 & 56.9 & \\
\hline
\end{tabular}
\end{center}
Here, the downwind side contains the majority of AC area but not the majority
of the comparison-grid population. Thus
\texttt{downup\_ac\_area}=1 and
\texttt{downup\_ac\_pop}=0.

\clearpage
\begin{figure}[p]
\centering
\includegraphics[width=\textwidth]{downup\_area\_population\_classification\_figure.pdf}
\caption{Corrected construction for one real grid-month. Panel (a) uses the
East-referenced wind vector and a perpendicular line through the focal-grid
centroid. Panel (b) intersects the AC with the two complementary squares.
Panel (c) classifies every other same-AC population centroid using the same
wind direction; marker size represents population. Panel (d) shows why the
two majority indicators can differ.}
\end{figure}
\clearpage

\section*{Step 1: construct the area measure}
Let $p=(x_i,y_i)$ be the focal-grid centroid and let $\theta$ be the rolling
wind direction from East. Define
\[
v=(\cos\theta,\sin\theta),\qquad
t=(\sin\theta,-\cos\theta),
\]
where $v$ points downwind and $t$ is perpendicular to the wind. Choose $R$
larger than the maximum distance from $p$ to the AC bounds. The downwind square
has corners
\[
p-Rt,\quad p+Rt,\quad p+Rt+2Rv,\quad p-Rt+2Rv.
\]
The upwind square uses $-v$. We compute the exact intersections
\[
A_{down}=\mathrm{area}(AC\cap S_{down}),\qquad
A_{up}=\mathrm{area}(AC\cap S_{up}).
\]
The binary indicator equals 1 only when $A_{down}>A_{up}$. An exact tie is
coded 0. A missing rolling direction produces a missing indicator.

\section*{Step 2: construct the population measure}
For every other grid $j$ assigned to the same AC, compute the signed circular
difference
\[
\delta_{ij}=((\beta_{ij}-\theta+180)\bmod 360)-180.
\]
Grid $j$ is downwind when $|\delta_{ij}|<90^\circ$ and upwind when
$|\delta_{ij}|>90^\circ$. The focal grid is excluded because its centroid is
on the dividing line. A comparison centroid exactly on that line is also
excluded from both sums. Population is then
\[
P_{down}=\sum_{j:\,|\delta_{ij}|<90^\circ} Pop_{j,2010},\qquad
P_{up}=\sum_{j:\,|\delta_{ij}|>90^\circ} Pop_{j,2010}.
\]
The binary indicator equals 1 only when $P_{down}>P_{up}$; ties are 0 and
missing rolling direction produces a missing value.

\section*{What is shared and what is different}
\begin{center}
\begin{tabular}{lll}
\hline
Component & Area construction & Population construction \\
\hline
Focal unit & grid centroid & grid centroid \\
Wind angle & rolling, from East & rolling, from East \\
Divider & perpendicular line & same perpendicular line \\
Objects classified & AC polygon surface & same-AC grid centroids \\
Weights & square kilometres & 2010 grid population \\
Indicator & down area $>$ up area & down pop. $>$ up pop. \\
\hline
\end{tabular}
\end{center}

\clearpage
\section*{Implementation safeguards}
The production code retains the unrounded rolling-direction float for every
grid-month. Area rows are never rounded, grouped, cached, or deduplicated by
wind direction. The population lookup uses circular intervals and an ASOF join
only after both the wind and centroid bearings have been placed in the same
East-referenced $[0,360)$ system.

The automated tests enforce the four cardinal cases:
\begin{itemize}
  \item East wind: the eastern half-plane is downwind.
  \item North wind: the northern half-plane is downwind.
  \item West wind: the western half-plane is downwind.
  \item South wind: the southern half-plane is downwind.
\end{itemize}
They also test the $179^\circ/-179^\circ$ wraparound to ensure that monthly and
rolling directions are calculated from eastward/northward vector components,
not arithmetic averages of degree values.

\section*{Production sequence}
The corrected files are generated in this order:
\begin{enumerate}
  \item \texttt{build\_wind\_direction\_grid\_duckdb.py}\par
  \item \texttt{build\_downup\_ac\_pop\_cluster.py}\par
  \item \texttt{build\_downup\_ac\_area\_cluster.py}\par
  \item \texttt{build\_downup\_13kmpl\_cluster.py}\par
  \item \texttt{build\_0\_master\_dataset.py}\par
  \item \texttt{build\_all\_stacked\_datasets\_duckdb.py}\par
\end{enumerate}
The one-job launcher is
\texttt{rebuild\_corrected\_downup\_pipeline.sbatch}.

\section*{Interpretation}
Both indicators describe exposure generated by a hypothetical fire at the
focal grid under the prevailing rolling wind. The area measure asks whether
most of the constituency's land lies downwind. The population measure asks
whether most of the measured same-AC population lies downwind. Neither measure
should be interpreted as the other, and disagreement between them is a valid
feature of the construction rather than a merge error.

\end{document}
